% Generated from tex/chapters/ch04/08_構築技術.tex for the SWEBOK structure tree
\subsubsection{構築のためのフィードバックループ}

構築で重要なのは、早く継続的なフィードバックを得ることである。アジャイル開発では、短い反復で動くソフトウェアを作り、顧客やチームから反応を得る。DevOps では、運用環境からの情報を開発へ戻すことで、実際の性能、信頼性、利用状況を学ぶ。

フィードバックループを支える技法には、自動ビルド、自動テスト、継続的統合、カナリアリリース、A/B テスト、監視、ログ分析がある。カナリアリリースは、全利用者に公開する前に、一部の利用者や環境へ変更を出して問題を観察する方法である。A/B テストは、複数の方式を比較し、実利用のデータから判断する方法である。
